\section*{第 \ref{ch:g4:adapted-to-their-world} 章——适应他们的世界}

\begin{solution}{exo:g4:adapted-to-their-world:1}
使\omterm{def:g1:living-or-not:living}{生物}\omterm{def:g4:adapted-to-their-world:adaptation}{适应}其\omterm{def:g2:caring-for-nature:environment}{环境}条件的身体或
行为特征。\omterm{def:g1:animals-around-us:animal}{动物}：北极熊毛皮与脂肪的双重保温；
\omterm{def:g1:plants-around-us:plant}{植物}：仙人掌储存水分的
茎。
\end{solution}

\begin{solution}{exo:g4:adapted-to-their-world:2}
缺水：一次能饮下一百升水，数日不饮仍可存活，
并靠驼峰中的脂肪储备维持。松软沙地：宽阔的
两趾足。风沙：可闭合的鼻孔和双层长睫毛。
（炎热与饥饿：又是驼峰的作用，脂肪
可供穿越贫瘠地带时使用。）
\end{solution}

\begin{solution}{exo:g4:adapted-to-their-world:3}
寒冷：厚密的毛皮之下是深厚的脂肪层。冰冷的水和积雪：
宽阔且略带蹼的爪——既像雪鞋，又像船桨。在白色冰面上
捕食警觉的猎物：白色\omterm{def:g1:animals-around-us:covering}{皮毛}的伪装作用。
\end{solution}

\begin{solution}{exo:g4:adapted-to-their-world:4}
仙人掌的刺是退化的叶片：几乎不散失水分，
还能使\omterm{def:g1:plants-around-us:plant}{植物}免受口渴\omterm{def:g1:animals-around-us:animal}{动物}的侵害。桶状的茎
储存水分，并承担了叶片捕获光能的
工作。
\end{solution}

\begin{solution}{exo:g4:adapted-to-their-world:5}
垫状的外形挡住了刺骨的寒风，使其内部的空气
保持相对温和。外形本身就是一种庇护。
\end{solution}

\begin{solution}{exo:g4:adapted-to-their-world:6}
作为猎物：保护色、盔甲或尖刺、速度，以及宽距的
眼睛可四面观察（具备两项即可）。作为捕食者：朝前的
眼睛、敏锐的感觉、爪和速度（具备两项即可）。
\end{solution}

\begin{solution}{exo:g4:adapted-to-their-world:7}
浓密柔软的\omterm{def:g1:animals-around-us:covering}{皮毛}——\omterm{prop:g3:sorting-living-things:groups}{哺乳动物}；巨大的铲状前爪；
——善于挖掘；细小而弱的眼睛——\omterm{def:g4:adapted-to-their-world:adaptation}{适应}黑暗生活。综合起来，
这是一种地下挖掘的\omterm{prop:g3:sorting-living-things:groups}{哺乳动物}：鼹鼠。
\end{solution}

\begin{solution}{exo:g4:adapted-to-their-world:8}
答案不唯一。以鸭子为例：蹼状足（便于划水），
防水\omterm{def:g1:animals-around-us:covering}{羽毛}（浮在水面时保持干燥和温暖），扁平的
滤食喙（从水中获取食物）。
\end{solution}

\begin{solution}{exo:g4:adapted-to-their-world:9}
耳朵是散热器：又大又薄的耳朵能散发体热。在
沙漠中，散热是需要解决的问题——巨大的
耳朵可以为狐狸降温。在冰原上，保温才是关键——
细小的耳朵散热最少。同一器官，两种气候，两种
截然相反的大小。
\end{solution}

\begin{solution}{exo:g4:adapted-to-their-world:10}
流线型身体、蹼状足和能屏住\omterm{def:g4:breathing:movements}{呼吸}，说明它是游泳者
和潜水者；致密防水的\omterm{def:g1:animals-around-us:covering}{皮毛}说明它是半生活在冷水中的
\omterm{prop:g3:sorting-living-things:groups}{哺乳动物}；胡须则说明它依靠触觉在黑暗或浑浊的水中探路。
这是一种生活在水中、会筑坝的\omterm{prop:g3:sorting-living-things:groups}{哺乳动物}，生活在河流
和湖泊中——海狸。
\end{solution}

\begin{solution}{exo:g4:adapted-to-their-world:11}
决定其形态的是水。鲸和鳟鱼并未从共同祖先那里遗传
到适于游泳的形态——一个是\omterm{prop:g3:sorting-living-things:groups}{哺乳动物}，另一个是\omterm{prop:g3:sorting-living-things:groups}{鱼类}，
但面对同样的问题，即在水中快速移动，两者都演化出
相同的流线型结构。\omterm{def:g4:adapted-to-their-world:adaptation}{适应}是对\omterm{def:g2:caring-for-nature:environment}{环境}问题的回应，
相似的问题可以在亲缘关系很远的类群中塑造出
相似的形态。
\end{solution}
